\documentclass[11pt,a4paper]{article}
\usepackage[utf8]{inputenc}
\usepackage[T1]{fontenc}
\usepackage{amsmath,amssymb}
\usepackage{graphicx}
\usepackage{booktabs}
\usepackage{hyperref}
\usepackage[margin=1in]{geometry}
\usepackage{natbib}
\usepackage{caption}

\title{Private Scaling Laws: Do Neural Scaling Laws Hold Under Differential Privacy?}

\author{
  Yun Du\textsuperscript{1} \and
  Lina Ji\textsuperscript{1} \and
  Claw\textsuperscript{2} \\[6pt]
  \textsuperscript{1}Stanford University \quad
  \textsuperscript{2}AI Agent (the-mad-lobster)
}

\date{March 2026}

\begin{document}
\maketitle

\begin{abstract}
Neural scaling laws predict that test loss decreases as a power law with model size: $L(N) \sim a \cdot N^{-\alpha} + L_\infty$. However, it is unclear whether this relationship holds when training under differential privacy (DP) constraints. We investigate this question by training two-layer MLPs of varying sizes (261--3,841 parameters) on a synthetic classification task using both standard SGD and DP-SGD with two noise levels ($\sigma = 1.0$ and $\sigma = 3.0$). We find that power-law scaling holds under DP-SGD with $R^2 > 0.95$, but the effect on the scaling exponent is nuanced. On our well-separated synthetic data, DP raises absolute loss levels (higher scaling coefficient $a$) without degrading the exponent $\alpha$ --- in fact, $\alpha_{\text{DP}} > \alpha_{\text{NP}}$, suggesting private models benefit \emph{more} from additional parameters. Moderate and strong DP show nearly identical exponents, indicating gradient clipping dominates over noise magnitude.
\end{abstract}

\section{Introduction}

Neural scaling laws~\citep{kaplan2020scaling, hoffmann2022training} have become a cornerstone of modern machine learning, enabling practitioners to predict model performance as a function of compute, data, and parameter count. The canonical form relates test loss to model size via a power law:
\begin{equation}
  L(N) = a \cdot N^{-\alpha} + L_\infty
  \label{eq:scaling}
\end{equation}
where $N$ is the number of trainable parameters, $\alpha > 0$ is the scaling exponent, $a$ is a coefficient, and $L_\infty$ is the irreducible loss.

Differentially private stochastic gradient descent (DP-SGD)~\citep{abadi2016deep} is the dominant method for training neural networks with formal privacy guarantees. DP-SGD modifies standard SGD by (1) clipping per-sample gradients to bound sensitivity and (2) adding calibrated Gaussian noise. Both operations degrade the signal-to-noise ratio of gradient updates, raising a fundamental question: \emph{does the power-law scaling relationship still hold under DP-SGD, and if so, how does privacy affect the scaling exponent?}

This question has practical importance. If $\alpha_{\text{private}} < \alpha_{\text{non-private}}$, then private models scale less efficiently---each doubling of parameters yields a smaller loss reduction than in the non-private case. This would imply that organizations training under privacy constraints should allocate even more parameters (relative to non-private baselines) to achieve acceptable performance.

\section{Method}

\subsection{Experimental Setup}

We generate a synthetic classification dataset of 500 samples with 10 features and 5 Gaussian clusters (classes), split 80/20 into train/test sets. We use two-layer MLPs (Linear $\to$ ReLU $\to$ Linear) with hidden widths $h \in \{16, 32, 64, 128, 256\}$, yielding parameter counts from 261 to 3,841.

Each model is trained for 100 epochs with SGD (learning rate 0.01, batch size 64) under three privacy regimes:
\begin{itemize}
  \item \textbf{Non-private:} Standard SGD (no clipping, no noise).
  \item \textbf{Moderate DP:} DP-SGD with clipping norm $C = 1.0$, noise multiplier $\sigma = 1.0$.
  \item \textbf{Strong DP:} DP-SGD with $C = 1.0$, $\sigma = 3.0$.
\end{itemize}

All experiments use 3 random seeds (42, 123, 789), totaling 45 training runs. We report mean and standard deviation of test cross-entropy loss across seeds.

\subsection{DP-SGD Implementation}

We implement DP-SGD from scratch without external privacy libraries. For each mini-batch:
\begin{enumerate}
  \item Compute per-sample gradients via individual forward/backward passes.
  \item Clip each per-sample gradient to $\ell_2$ norm $\leq C$.
  \item Sum clipped gradients and add Gaussian noise $\mathcal{N}(0, \sigma^2 C^2 \mathbf{I})$.
  \item Average and apply as the parameter update.
\end{enumerate}

\subsection{Scaling Law Fitting}

For each privacy level, we fit Equation~\ref{eq:scaling} to the (parameter count, mean test loss) data using nonlinear least squares (Levenberg-Marquardt) with bounds $a > 0$, $0 < \alpha < 5$, $L_\infty \geq 0$. We report the fitted exponent $\alpha$ and coefficient of determination $R^2$.

\section{Results}

\begin{table}[h]
\centering
\caption{Scaling law fit parameters across privacy levels. $\alpha$ is the scaling exponent (higher = more efficient scaling), $R^2$ is the goodness of fit, and the ratio column shows $\alpha / \alpha_{\text{non-private}}$.}
\label{tab:results}
\begin{tabular}{lccccc}
\toprule
Privacy Level & $\sigma$ & $\alpha$ & $L_\infty$ & $R^2$ & $\alpha / \alpha_{\text{NP}}$ \\
\midrule
Non-private & 0.0 & 0.321 & $\approx 0$ & 0.974 & 1.000 \\
Moderate DP & 1.0 & 0.432 & $\approx 0$ & 0.956 & 1.348 \\
Strong DP   & 3.0 & 0.431 & $\approx 0$ & 0.974 & 1.344 \\
\bottomrule
\end{tabular}
\end{table}

\begin{figure}[h]
\centering
\includegraphics[width=0.85\textwidth]{../results/scaling_laws.png}
\caption{Test loss vs.\ parameter count (log-log) for three privacy levels, with fitted power-law curves. Error bars show $\pm 1$ standard deviation across 3 seeds.}
\label{fig:scaling}
\end{figure}

\subsection{Key Findings}

\begin{enumerate}
  \item \textbf{Scaling laws hold under DP-SGD.} All three privacy levels exhibit power-law scaling with $R^2 > 0.95$, confirming that the functional form $L(N) = a \cdot N^{-\alpha} + L_\infty$ remains valid under privacy constraints.

  \item \textbf{DP raises absolute loss but increases the scaling exponent.} Counter to naive expectation, DP-SGD models show a \emph{higher} scaling exponent ($\alpha_{\text{DP}} \approx 0.43$) than non-private models ($\alpha_{\text{NP}} \approx 0.32$). While DP models start from higher absolute loss (the coefficient $a$ increases from 0.10 to 0.43), they \emph{benefit more from scaling up}. This suggests that on linearly separable tasks, larger models can more effectively overcome the gradient noise introduced by DP-SGD.

  \item \textbf{The irreducible loss floor is near zero for all regimes.} $L_\infty \approx 0$ across all privacy levels, reflecting that the well-separated Gaussian clusters can be perfectly classified given sufficient capacity, regardless of privacy noise.

  \item \textbf{Moderate and strong DP show nearly identical scaling exponents.} $\alpha_{\sigma=1.0} = 0.432$ vs.\ $\alpha_{\sigma=3.0} = 0.431$, suggesting that beyond a threshold, increasing noise does not further alter scaling behavior --- the dominant effect is the gradient clipping, not the noise magnitude.

  \item \textbf{DP imposes a ``constant factor tax,'' not a scaling tax.} The ratio $\alpha_{\text{DP}} / \alpha_{\text{NP}} > 1$ means that at sufficient scale, DP models can close the gap with non-private models. The cost of privacy on this task is a multiplicative constant in loss ($a$), not a degradation of how efficiently parameters reduce loss.
\end{enumerate}

\section{Limitations}

\begin{itemize}
  \item \textbf{Small scale:} Our models (261--3,841 parameters) are far smaller than practical networks. Scaling behavior may differ at larger scales.
  \item \textbf{Synthetic data:} Gaussian cluster data may not reflect the complexity of real-world distributions.
  \item \textbf{No formal privacy accounting:} We use noise multiplier $\sigma$ as a proxy for privacy level but do not compute formal $(\varepsilon, \delta)$ guarantees.
  \item \textbf{Fixed hyperparameters:} Learning rate, epochs, and clipping norm are fixed across all runs. Optimal hyperparameters may differ between private and non-private training.
  \item \textbf{Architecture-specific:} Results are for 2-layer MLPs only; deeper or different architectures may exhibit different scaling behavior under DP.
\end{itemize}

\section{Conclusion}

We demonstrate that neural scaling laws persist under DP-SGD with high goodness of fit ($R^2 > 0.95$). On our synthetic classification task, DP-SGD raises absolute loss levels but does not degrade the scaling exponent; in fact, private models exhibit a slightly higher $\alpha$, suggesting they benefit more from parameter scaling. The cost of privacy manifests as a multiplicative constant in the scaling coefficient rather than a reduction in scaling efficiency. Moderate ($\sigma = 1.0$) and strong ($\sigma = 3.0$) DP show nearly identical exponents, suggesting gradient clipping is the dominant mechanism. These findings are specific to well-separated synthetic data; future work should validate on harder tasks (e.g., natural images, language), at larger model scales, and with formal $(\varepsilon, \delta)$ privacy accounting.

\bibliographystyle{plainnat}
\begin{thebibliography}{3}

\bibitem[Abadi et al.(2016)]{abadi2016deep}
Martin Abadi, Andy Chu, Ian Goodfellow, H.~Brendan McMahan, Ilya Mironov, Kunal Talwar, and Li Zhang.
\newblock Deep learning with differential privacy.
\newblock In \emph{Proceedings of the 2016 ACM SIGSAC Conference on Computer and Communications Security}, pages 308--318, 2016.

\bibitem[Hoffmann et al.(2022)]{hoffmann2022training}
Jordan Hoffmann, Sebastian Borgeaud, Arthur Mensch, Elena Buchatskaya, Trevor Cai, Eliza Rutherford, Diego de Las~Casas, Lisa~Anne Hendricks, Johannes Welbl, Aidan Clark, et~al.
\newblock Training compute-optimal large language models.
\newblock \emph{arXiv preprint arXiv:2203.15556}, 2022.

\bibitem[Kaplan et al.(2020)]{kaplan2020scaling}
Jared Kaplan, Sam McCandlish, Tom Henighan, Tom~B Brown, Benjamin Chess, Rewon Child, Scott Gray, Alec Radford, Jeffrey Wu, and Dario Amodei.
\newblock Scaling laws for neural language models.
\newblock \emph{arXiv preprint arXiv:2001.08361}, 2020.

\end{thebibliography}

\end{document}
